\documentclass[11pt,a4paper]{article}

\usepackage[T1]{fontenc}
\usepackage[utf8]{inputenc}
\usepackage{mathpazo}
\usepackage{microtype}
\usepackage[a4paper,margin=27mm,headheight=23pt]{geometry}
\usepackage{amsmath,amssymb,amsthm,mathtools}
\usepackage{booktabs,tabularx,array,longtable}
\usepackage{enumitem}
\usepackage{xcolor}
\usepackage{fancyhdr}
\usepackage{lastpage}
\usepackage{hyperref}
\usepackage{xurl}
\usepackage[nameinlink,noabbrev]{cleveref}

\definecolor{ink}{HTML}{172033}
\definecolor{accent}{HTML}{245A86}
\definecolor{accentlight}{HTML}{EAF2F8}
\definecolor{warning}{HTML}{8A4B08}
\definecolor{warninglight}{HTML}{FFF5E6}
\definecolor{rulegray}{HTML}{CBD5E1}

\hypersetup{
  colorlinks=true,
  linkcolor=accent,
  citecolor=accent,
  urlcolor=accent,
  pdfauthor={Prepared from the z20 reproduction bundle},
  pdftitle={A Candidate Certificate-Backed Proof that z(20)=6},
  pdfsubject={Candidate computer-assisted proof of Erdos Problem 758},
  pdfkeywords={cochromatic number, Ramsey graph, SAT certificate, DRUP, LRAT}
}
\urlstyle{same}

\pagestyle{fancy}
\fancyhf{}
\lhead{\small\textcolor{ink}{A candidate certificate-backed proof that \(z(20)=6\)}}
\rhead{\small\textcolor{ink}{Candidate proof}}
\cfoot{\small\textcolor{ink}{\thepage\ of \pageref*{LastPage}}}
\renewcommand{\headrulewidth}{0.4pt}
\renewcommand{\headrule}{\hbox to\headwidth{\color{rulegray}\leaders\hrule height \headrulewidth\hfill}}

\setlength{\parindent}{0pt}
\setlength{\parskip}{0.62em}
\setlist[itemize]{leftmargin=1.35em,itemsep=0.25em,topsep=0.35em}
\setlist[enumerate]{leftmargin=1.6em,itemsep=0.3em,topsep=0.35em}
\renewcommand{\arraystretch}{1.17}

\newtheorem{theorem}{Theorem}[section]
\newtheorem{lemma}[theorem]{Lemma}
\newtheorem{proposition}[theorem]{Proposition}
\newtheorem{corollary}[theorem]{Corollary}
\theoremstyle{definition}
\newtheorem{definition}[theorem]{Definition}
\theoremstyle{remark}
\newtheorem{remark}[theorem]{Remark}

\newcommand{\zetaG}{\zeta}
\newcommand{\homfour}{\textnormal{hom-}4}
\newcommand{\corezero}{\texttt{\detokenize{OsHHirKdlp[IFVI|KpqfR}}}
\newcommand{\coreone}{\texttt{\detokenize{Ov?IXZIhlRWjUXL[iphst}}}
\newcommand{\sha}[1]{\texttt{\footnotesize #1}}
\newcommand{\artifact}[1]{\texttt{\detokenize{#1}}}
\newcommand{\statusbox}[1]{%
  \begin{center}
  \fcolorbox{warning}{warninglight}{%
    \parbox{0.88\textwidth}{\small\textcolor{warning}{\textbf{Status.}} #1}}
  \end{center}}
\newcommand{\resultbox}[1]{%
  \begin{center}
  \fcolorbox{accent}{accentlight}{%
    \parbox{0.88\textwidth}{\centering\large\textcolor{accent}{#1}}}
  \end{center}}

\begin{document}
\hypersetup{pageanchor=false}

\begin{titlepage}
  \thispagestyle{empty}
  \vspace*{2.1cm}
  {\color{accent}\rule{\textwidth}{1.2pt}}
  \vspace{1.1cm}

  {\Huge\bfseries\color{ink} A Candidate Certificate-Backed Proof\\[0.18em]
  that \(z(20)=6\)\par}

  \vspace{0.65cm}
  {\Large\color{accent} The cochromatic number of graphs on twenty vertices\\
  (Erd\H{o}s Problem \#758)\par}

  \vspace{1.25cm}
  {\large Prepared from the local \texttt{z20\_repro} verification bundle\par}
  \vspace{0.25cm}
  {\large 26 July 2026\par}

  \vfill

  \statusbox{This is an \textbf{unrefereed candidate computer-assisted proof},
  offered for independent mathematical and computational scrutiny.  The local
  certificates have been replayed and the feasible exhaustive computations
  have been rerun.  No human peer review, external reproduction, or
  end-to-end proof-assistant formalisation is claimed.}

  \vfill
  {\color{accent}\rule{\textwidth}{1.2pt}}
\end{titlepage}
\hypersetup{pageanchor=true}
\pagenumbering{arabic}

\begin{abstract}
For a finite simple graph \(G\), let \(\zeta(G)\) be the least number of
classes in a partition of \(V(G)\) such that each class induces either a
clique or an independent set, and put
\[
  z(n)=\max\{\zeta(G):|V(G)|=n\}.
\]
The next value listed as unresolved is \(z(20)\), known to be either \(6\)
or \(7\).  This note presents a candidate computer-assisted proof that
\(z(20)=6\).

The lower bound is a hand proof using the Paley graph on \(17\) vertices.
For the upper bound, every \(20\)-vertex graph is shown to contain two
vertex-disjoint homogeneous \(4\)-sets.  If one fixes a first such set, a
putative counterexample leaves a homogeneous-\(4\)-free graph on \(16\)
vertices.  Orderly generation gives exactly two such graphs up to
isomorphism.  Each leaves \(64\) undetermined cross edges and yields a
\(64\)-variable CNF with \(104{,}524\) clauses.  Both formulas have
certificate-backed refutations.  Removing the two homogeneous \(4\)-sets
leaves \(12\) vertices, and a separate exhaustive computation proves
\(z(12)\leq4\).

The note separates the mathematical reduction, the exact encoding lemma,
the proof-certificate checks, and the remaining trusted-computing
assumptions.  In particular, the SAT checkers certify the two CNFs, not the
graph-theoretic theorem by themselves.
\end{abstract}

\noindent\textbf{Keywords:} cochromatic number; Ramsey graph; exhaustive
generation;\\
SAT certificate; DRUP; LRAT.

\tableofcontents
\newpage

\section{Statement, scope, and proof architecture}

Cocolouring was introduced by Lesniak-Foster and Straight
\cite{lesniak-straight}.  The finite extremal question considered here is
recorded as Erd\H{o}s Problem \#758 \cite{erdos758}.  This document is based
on the local repository \artifact{/Users/admin/z20_repro} and the public
artifact snapshot \cite{z20repo}.

\begin{definition}
A set \(S\subseteq V(G)\) is \emph{homogeneous} if \(G[S]\) is complete or
edgeless.  A \emph{cocolouring} of \(G\) is a partition of \(V(G)\) into
homogeneous classes.  Its minimum number of classes is the
\emph{cochromatic number} \(\zeta(G)\).  Finally,
\[
  z(n)=\max\{\zeta(G):G\text{ is a finite simple graph on }n\text{ vertices}\}.
\]
A \((4,4)\)-graph is a graph with neither a \(K_4\) nor an independent
\(4\)-set.
\end{definition}

\begin{theorem}[Main result]\label{thm:main}
\[
  \boxed{z(20)=6}.
\]
\end{theorem}

The proof has two unequal halves.  The lower bound is entirely
hand-checkable.  The upper bound depends on three finite computations:

\begin{center}
\begin{tabularx}{0.94\textwidth}{@{}>{\bfseries}l X@{}}
\toprule
Obligation & Role in the proof\\
\midrule
Catalogue &
There are exactly two \((4,4)\)-graphs on \(16\) vertices, up to
isomorphism; the same run supplies a \(17\)-vertex survivor and no
\(18\)-vertex survivor.\\
Fixed-core SAT &
For neither \(16\)-vertex core can the \(64\) cross edges be chosen so as
to avoid two disjoint homogeneous \(4\)-sets.\\
Residual cocolouring &
Every \(8\)-vertex graph is \(3\)-cocolourable, and every
\((4,4)\)-graph on \(12\) vertices is \(4\)-cocolourable.\\
\bottomrule
\end{tabularx}
\end{center}

The logical flow is
\[
\begin{gathered}
  P_{17}\text{ has }\alpha(P_{17}),\omega(P_{17})\leq3
      \quad\Longrightarrow\quad z(20)\geq6,\\[0.35em]
  R(4,4)=18+\text{two }16\text{-vertex cores}
      +\text{two certified UNSAT results}\\
  \Longrightarrow\text{ two disjoint homogeneous }4\text{-sets in every }G_{20}\\
  \Longrightarrow 2+z(12)\leq 2+4=6
      \quad\Longrightarrow\quad z(20)\leq6.
\end{gathered}
\]

\section{Elementary invariances}

\begin{lemma}[Induced-subgraph monotonicity]\label{lem:monotone}
If \(H\) is an induced subgraph of \(G\), then
\(\zeta(H)\leq\zeta(G)\).
\end{lemma}

\begin{proof}
Restrict a cocolouring of \(G\) to \(V(H)\).  Every resulting class is a
subset of a clique or of an independent set, hence is still homogeneous.
Empty classes may be discarded.
\end{proof}

\begin{lemma}[Complement invariance]\label{lem:complement}
A vertex set is homogeneous in \(G\) if and only if it is homogeneous in
\(\overline G\).  Consequently, the property
\[
  \text{``has two vertex-disjoint homogeneous \(4\)-sets''}
\]
is invariant under complementation.
\end{lemma}

\begin{proof}
Complementation exchanges cliques and independent sets and leaves the
underlying vertex sets unchanged.
\end{proof}

\section{The lower bound: a hand proof}

Let
\[
  Q=\{1,2,4,8,9,13,15,16\}\subset\mathbb Z_{17}
\]
be the set of nonzero quadratic residues.  The Paley graph \(P_{17}\) has
vertex set \(\mathbb Z_{17}\), with \(i\sim j\) precisely when
\(i-j\in Q\).  Since \(-1=16\in Q\), this defines an undirected graph.

\begin{lemma}\label{lem:paley}
The Paley graph \(P_{17}\) contains neither a \(K_4\) nor an independent
\(4\)-set.
\end{lemma}

\begin{proof}
Translations and multiplications by a nonzero square are automorphisms of
\(P_{17}\).  They act transitively on edges: after translating one endpoint
to \(0\), multiplication by the inverse of the edge difference sends the
other endpoint to \(1\).  Hence a hypothetical \(K_4\) may be assumed to
contain the edge \(\{0,1\}\).

The common neighbours of \(0\) and \(1\) are
\[
  Q\cap(1+Q)=\{2,9,16\}.
\]
Their pairwise differences are congruent, up to sign, to \(7\) or \(14\);
neither lies in \(Q\).  Thus the three common neighbours are pairwise
nonadjacent, so no two of them can complete \(\{0,1\}\) to a \(K_4\).

Multiplication by any quadratic nonresidue, for example \(3\), sends
residues to nonresidues and is an isomorphism
\(P_{17}\cong\overline{P_{17}}\).  The absence of a \(K_4\) therefore also
implies the absence of an independent \(4\)-set.
\end{proof}

\begin{proposition}\label{prop:lower}
\(z(20)\geq6\).
\end{proposition}

\begin{proof}
By \cref{lem:paley}, every homogeneous set in \(P_{17}\) has at most three
vertices.  Five homogeneous classes could therefore cover at most
\(15\) vertices, so
\[
  \zeta(P_{17})\geq \left\lceil\frac{17}{3}\right\rceil=6.
\]
Let \(G=P_{17}\sqcup3K_1\).  The graph \(G\) has \(20\) vertices and contains
\(P_{17}\) as an induced subgraph.  By \cref{lem:monotone},
\(\zeta(G)\geq6\), proving \(z(20)\geq6\).
\end{proof}

For completeness, the following partition is an explicit six-class
cocolouring of \(G\), where \(17,18,19\) are the isolated vertices:
\[
\begin{aligned}
 I_1&=\{0,3,6,17,18,19\},&
 I_2&=\{1,4,7\},&
 I_3&=\{2,9,12\},\\
 I_4&=\{5,8,11\},&
 I_5&=\{10,13,16\},&
 K_1&=\{14,15\}.
\end{aligned}
\]
The \(I_i\) are independent sets and \(K_1\) is an edge.  No solver is
needed for the lower bound.

\section{\texorpdfstring{Two disjoint homogeneous \(4\)-sets}{Two disjoint homogeneous 4-sets}}

The classical identity \(R(4,4)=18\) is due to Greenwood and Gleason
\cite{greenwood-gleason}; it is also recomputed by the catalogue run
described in \cref{sec:catalogue}.

\begin{lemma}\label{lem:two-homfour}
Every graph on \(20\) vertices contains two vertex-disjoint homogeneous
\(4\)-sets.
\end{lemma}

\begin{proof}
Suppose for contradiction that a \(20\)-vertex graph \(G\) has no two
vertex-disjoint homogeneous \(4\)-sets.  Since \(20\geq R(4,4)=18\), it contains a
homogeneous \(4\)-set \(A\).  By \cref{lem:complement}, complementing \(G\)
if needed does not affect the target property, so assume that \(A\) is a
clique.  Relabel \(A=\{a_0,a_1,a_2,a_3\}\) and put
\(B=V(G)\setminus A\).

If \(G[B]\) contains a homogeneous \(4\)-set, it is disjoint from \(A\) and
the conclusion follows.  Otherwise \(G[B]\) is a \((4,4)\)-graph on
\(16\) vertices.  The complete catalogue in
\cref{sec:catalogue} gives, up to isomorphism, exactly the following two
cores in graph6 notation:
\[
\begin{array}{ll}
C_0=&\corezero,\\
C_1=&\coreone.
\end{array}
\]
Each has \(60\) edges, degree sequence \(7^8 8^8\), \(56\) triangles, and
\(56\) independent triples.

After fixing \(G[B]\) to \(C_0\) or \(C_1\), the only undetermined edges are
the \(4\cdot16=64\) edges between \(A\) and \(B\).  For each core,
\cref{sec:encoding} constructs a CNF whose satisfying assignments are
\emph{exactly} those cross-edge patterns for which the resulting graph has
no two disjoint homogeneous \(4\)-sets.  The actual cross-edge pattern of
\(G\) would therefore satisfy one of these CNFs, but both are unsatisfiable
by the certificates in \cref{sec:certificates}.  This contradiction proves
the claim.

Notice that \(A\) itself is a homogeneous \(4\)-set.  It need not be
encoded.  The set \(B\) contains no homogeneous \(4\)-set, so \(A\) cannot
be one half of a disjoint pair.  Every other possible pair consists of two
mixed sets, which the CNF treats exhaustively.
\end{proof}

\section{The exact fixed-core encoding}\label{sec:encoding}

Fix \(A=\{a_0,a_1,a_2,a_3\}\) as a clique and let
\(B=\{b_0,\ldots,b_{15}\}\) induce one of the two fixed cores.  For
\(0\leq i<4\) and \(0\leq j<16\), introduce
\[
  x_{i,j}=\begin{cases}
  1,&a_i b_j\in E(G),\\
  0,&a_i b_j\notin E(G).
  \end{cases}
\]

\begin{lemma}[Mixed-set classification]\label{lem:mixed}
The mixed \(4\)-sets that can be homogeneous are exactly the four types in
\cref{tab:mixed}.
\end{lemma}

\begin{table}[ht]
\centering
\caption{Conditions under which a mixed \(4\)-set is homogeneous.}
\label{tab:mixed}
\begin{tabularx}{\textwidth}{@{}l X r@{}}
\toprule
Split & Necessary and sufficient condition & Count\\
\midrule
\(3A+1B\) &
Clique: the three relevant cross edges are present. &
\(\binom43\cdot16=64\)\\
\(2A+2B\) &
Clique: the \(B\)-pair is one of the \(60\) core edges and all four
relevant cross edges are present. &
\(\binom42\cdot60=360\)\\
\(1A+3B\) &
Clique: the \(B\)-triple is one of the \(56\) triangles and all three
cross edges are present. &
\(4\cdot56=224\)\\
\(1A+3B\) &
Independent: the \(B\)-triple is one of the \(56\) independent triples
and all three cross edges are absent. &
\(4\cdot56=224\)\\
\midrule
&\textbf{Total}&\(\mathbf{872}\)\\
\bottomrule
\end{tabularx}
\end{table}

\begin{proof}
If a \(4\)-set contains three vertices of \(A\), those vertices already
form a triangle; the set can only be a clique, and this occurs precisely
when the three cross edges are present.  If it contains two vertices of
\(A\), their edge rules out independence.  The set is a clique precisely
when the \(B\)-pair is an edge and the four cross edges are present.

With one vertex of \(A\), the \(B\)-triple must be either a triangle,
followed by three present cross edges, or an independent triple, followed
by three absent cross edges.  A \(4\)-set wholly in \(B\) is not
homogeneous by the definition of the core.  The only \(4\)-set wholly in
\(A\) is \(A\) itself, already handled in the last paragraph of
\cref{lem:two-homfour}.  This proves both necessity and sufficiency.
\end{proof}

Represent each possible mixed homogeneous set \(S\) by the conjunction
\(C_S\) of its three or four required cross-edge literals.  Whenever
\(S\cap T=\varnothing\), add the clause
\[
  \neg C_S\;\lor\;\neg C_T.
\]
Since \(C_S\) and \(C_T\) are conjunctions, this is a single clause
containing the negations of all their literals.  For example, if
\[
S=\{a_0,a_1,a_2,b_5\}
\quad\text{and}\quad
T=\{a_3,b_1,b_4,b_9\},
\]
where \(\{b_1,b_4,b_9\}\) is a core independent triple, the clause is
\[
\neg x_{0,5}\vee\neg x_{1,5}\vee\neg x_{2,5}\vee
x_{3,1}\vee x_{3,4}\vee x_{3,9}.
\]

\begin{proposition}[Encoding equivalence]\label{prop:encoding}
For either core, an assignment to the \(64\) variables satisfies the
constructed CNF if and only if the corresponding \(20\)-vertex graph has
no two vertex-disjoint homogeneous \(4\)-sets.
\end{proposition}

\begin{proof}
By \cref{lem:mixed}, every candidate other than \(A\) is represented by
exactly one condition \(C_S\), and the condition holds exactly when \(S\)
is homogeneous.  The clause for a disjoint pair \(S,T\) is false exactly
when both \(C_S\) and \(C_T\) hold.  Conjoining these clauses over all
disjoint pairs therefore forbids exactly, and only, the existence of two
disjoint homogeneous \(4\)-sets.
\end{proof}

For each core, the resulting DIMACS formula has \(64\) variables and
\(104{,}524\) clauses.  The clause-length histogram is
\[
  43{,}840\text{ of length }6,\qquad
  52{,}416\text{ of length }7,\qquad
  8{,}268\text{ of length }8.
\]
The only possible lengths are \(6,7,8\), the sums of two three- or
four-literal conditions subject to disjointness.

\section{Catalogue generation}\label{sec:catalogue}

McKay's maintained Ramsey catalogue lists the number of non-isomorphic
\((4,4)\)-graphs at each order and gives two at order \(16\)
\cite{mckay-catalogue}.  The local proof does not rely only on downloading
those two graphs: \artifact{nauty/gen44.c} regenerates the sequence from
the definition, using nauty for canonical labelling
\cite{mckay-piperno}.

\subsection{Why the generation is exhaustive}

Suppose a complete set of isomorphism-class representatives is known at
order \(n\).  Every \((4,4)\)-graph \(H\) on \(n+1\) vertices becomes a
\((4,4)\)-graph after deleting any vertex \(v\).  Hence \(H\) arises from
some order-\(n\) representative by adjoining one vertex with some
neighbourhood \(S\).  Trying every admissible \(S\), canonicalising the
result, and retaining one copy of each canonical form therefore reaches
every order-\(n+1\) isomorphism class.  Induction starts with the unique
one-vertex graph.

For a new vertex with neighbourhood \(S\), a new homogeneous \(4\)-set
must contain that vertex.  Admissibility is therefore equivalent to
\[
  G[S]\text{ contains no triangle}
  \quad\text{and}\quad
  G[V\setminus S]\text{ contains no independent triple}.
\]

The implementation prunes to \(n-8\leq|S|\leq8\).  Here is a short proof
of the needed bound, included because it is otherwise a hidden
\(R(3,4)\leq9\) input.  In a triangle-free graph on nine vertices with
independence number at most three, every degree is at most three, since
the neighbours of any vertex are independent.  If some vertex \(v\) had
degree at most two, deleting \(v\) and its neighbours would leave at
least six vertices.  The elementary \(R(3,3)\leq6\) argument gives an
independent triple among those remaining vertices, and adjoining \(v\)
would give an independent \(4\)-set.  Thus every degree would have to be
three, impossible on nine vertices by the handshaking lemma.  Therefore
a triangle-free set with independence number at most three has size at
most eight.  For completeness, the six-vertex argument is this: at any
vertex of a red--blue \(K_6\), three incident edges share a colour; their
other endpoints either contain an edge of that colour or form a triangle
of the other colour.  Thus \(R(3,3)\leq6\).

For \(V\setminus S\), the second admissibility condition says that
\(\overline{G[V\setminus S]}\) is triangle-free, while the absence of a
\(K_4\) in the core says that this complement has independence number at
most three.  The same bound gives \(|V\setminus S|\leq8\), hence
\(|S|\geq n-8\).

\subsection{Replayed output}

The full computation was replayed on 26 July 2026.  The counts, always
up to isomorphism, were:
\[
\begin{array}{c|rrrrrrrrr}
n&1&2&3&4&5&6&7&8&9\\ \hline
\#&1&2&4&9&24&84&362&2{,}079&14{,}701
\end{array}
\]
\[
\begin{array}{c|rrrrrrrrr}
n&10&11&12&13&14&15&16&17&18\\ \hline
\#&103{,}706&546{,}356&1{,}449{,}166&1{,}184{,}231&
130{,}816&640&2&1&0
\end{array}
\]
The two order-\(16\) graph6 strings are exactly \(C_0,C_1\) above.  The
order-\(17\) survivor supplies \(R(4,4)>17\), while the empty order-\(18\)
level supplies \(R(4,4)\leq18\); together they re-establish
\(R(4,4)=18\).

\section{Certificate verification}\label{sec:certificates}

For each fixed core, a clause-learning-free DPLL program produced a
tree-like DRUP proof.  Core \(0\) visited \(1{,}571\) search nodes, made
\(785\) decisions, and emitted \(1{,}571\) proof lines; core \(1\) visited
\(789\) nodes, made \(394\) decisions, and emitted \(789\) proof lines.
The preserved formulas have hashes
\[
\begin{array}{ll}
C_0:&
\sha{cbfea7b0b2cb712ad8b4f8b129f74c7f010e879064414eaef76e74aa88be3f44},\\[0.25em]
C_1:&
\sha{d3d96a3c2a719c753242673b81aabc3977c5e435bb5996b4c58c85a82d11dc5a}.
\end{array}
\]

The following checks were replayed on 26 July 2026.

\begin{longtable}{@{}>{\raggedright\arraybackslash}p{0.28\textwidth}
>{\raggedright\arraybackslash}p{0.31\textwidth}
>{\raggedright\arraybackslash}p{0.31\textwidth}@{}}
\caption{Replayed validation of the two fixed-core refutations.}
\label{tab:verification}\\
\toprule
Check & Core \(0\) & Core \(1\)\\
\midrule
\endfirsthead
\toprule
Check & Core \(0\) & Core \(1\)\\
\midrule
\endhead
Separately implemented formula construction &
\(872\) conditions; \(104{,}524\) clauses; clause set identical &
\(872\) conditions; \(104{,}524\) clauses; clause set identical\\
Direct semantic spot-check &
\(300/300\) assignments agreed &
\(300/300\) assignments agreed\\
Separately implemented RUP checker &
\(1{,}571\) steps; empty clause derived &
\(789\) steps; empty clause derived\\
Package RUP checker &
\texttt{VERIFIED} &
\texttt{VERIFIED}\\
\texttt{drat-trim} \cite{drat-trim} &
\texttt{s VERIFIED}; \(0\) RAT lemmas &
\texttt{s VERIFIED}; \(0\) RAT lemmas\\
\texttt{cake\_lpr} \cite{cake-lpr} &
\texttt{s VERIFIED UNSAT} &
\texttt{s VERIFIED UNSAT}\\
\bottomrule
\end{longtable}

The separately regenerated formulas are byte-identical to the
certificate-package formulas.  The package audit additionally tested
\(1{,}000\) deterministic random assignments per core.  The package's
twelve-entry SHA-256 manifest passes in full.

The DRUP proofs were converted to LRAT with \texttt{drat-trim -L} and the
resulting LRAT proofs were checked against the original CNFs by
\texttt{cake\_lpr}.  Once the LRAT proof is replayed against the CNF, the
translator need not be trusted for logical soundness: a faulty translation
would fail unless it happened to output another valid proof.

As a negative control, cross-pairing core \(0\)'s formula with core \(1\)'s
proof was rejected by the standalone RUP checker, by \texttt{drat-trim},
and by \texttt{cake\_lpr}.  One operational detail matters:
\texttt{cake\_lpr} can print a failure while returning process status \(0\).
Automation must require the exact success verdict
\texttt{s VERIFIED UNSAT}, not merely a zero exit status.

\section{\texorpdfstring{The residual \(12\)-vertex graph}{The residual 12-vertex graph}}

\begin{lemma}\label{lem:z8}
\(z(8)\leq3\).
\end{lemma}

\begin{proof}[Computational proof]
The program \texttt{geng} generated one representative of each of the
\(12{,}346\) isomorphism classes of simple graphs on eight vertices.
The exact recursive program \path{nauty/cochk.c} found every
representative \(3\)-cocolourable.
\end{proof}

\begin{lemma}\label{lem:z12}
\(z(12)\leq4\).
\end{lemma}

\begin{proof}
Let \(G\) have \(12\) vertices.  If \(G\) contains a homogeneous
\(4\)-set, use it as one class; by \cref{lem:z8}, the remaining eight
vertices require at most three classes.

Otherwise \(G\) is a \((4,4)\)-graph.  The orderly generator produced one
representative of every one of the \(1{,}449{,}166\) isomorphism classes
of such graphs on \(12\) vertices, and \path{nauty/cochk.c} found
each representative \(4\)-cocolourable.  These two cases are exhaustive.
\end{proof}

The recursive decision procedure is exact at the algorithmic level: choose
the least remaining vertex \(v\), enumerate every clique and every
independent set containing \(v\), remove the chosen block, and recurse with
one fewer available class.  Every cocolouring has exactly one class
containing \(v\), so no partition is omitted.  The assurance caveat is
implementation correctness, not an intentionally incomplete search.

\section{Completion of the upper bound}

\begin{proposition}\label{prop:upper}
\(z(20)\leq6\).
\end{proposition}

\begin{proof}
Let \(G\) have \(20\) vertices.  By \cref{lem:two-homfour}, choose two
disjoint homogeneous \(4\)-sets and use them as two cocolouring classes.
The remaining \(12\) vertices induce a graph cocolourable in at most four
classes by \cref{lem:z12}.  Thus \(\zeta(G)\leq2+4=6\).  Since \(G\) was
arbitrary, \(z(20)\leq6\).
\end{proof}

\begin{proof}[Proof of \cref{thm:main}]
\Cref{prop:lower} gives \(z(20)\geq6\), while \cref{prop:upper} gives
\(z(20)\leq6\).
\end{proof}

\resultbox{\(\displaystyle \boxed{z(20)=6}\)}

\section{Assurance boundary}\label{sec:assurance}

This section is part of the proof report, not a disclaimer to be read
after the fact.  Different components carry different assurance.

\begin{center}
\begin{tabularx}{\textwidth}{@{}>{\bfseries}p{0.22\textwidth} X@{}}
\toprule
Layer & What is established, and what remains trusted\\
\midrule
Hand mathematics &
The invariances, Paley lower bound, catalogue-completeness induction,
mixed-set classification, CNF equivalence, and final deductions can be
checked directly from this text.\\
SAT proof objects &
Each DRUP refutation is accepted by two separately implemented RUP
checkers and by \texttt{drat-trim}; the corresponding LRAT proof is
accepted against the exact CNF by the formally verified
\texttt{cake\_lpr} core.  This establishes unsatisfiability of those
CNFs.\\
Semantic bridge &
The graph-to-CNF equivalence is proved in \cref{prop:encoding} and is
also tested by a separately written direct graph implementation.
The source implementation of that equivalence is not formally verified.\\
Catalogue &
The generation is exhaustive if its implementation and nauty
canonicalisation are correct.  The complete sequence was replayed and
matches McKay's maintained data, but the C program and nauty are not
proved correct inside a proof assistant.\\
Cocolourability &
The recursive algorithm is exhaustive, and the full \(8\)- and
\(12\)-vertex runs were replayed.  The C implementation is not formally
verified.\\
System &
The compiler, operating system, file I/O, and hardware remain in the
conventional trusted computing base, except where the CakeML checker
reduces that base for the SAT-checking layer.\\
\bottomrule
\end{tabularx}
\end{center}

Two artifact-specific caveats should be visible to reproducers.

\begin{itemize}
  \item The recorded SHA-256 of \artifact{cake_lpr_arm8.S} matches the
  preserved assembly source.  The current \artifact{basis_ffi.c} hash does
  not match the older recorded hash.  The repository's explanation that
  this is a stale manifest for a small I/O shim is plausible, but is not
  established solely by the local files.
  \item \artifact{verify_certificate.py} contains a hard-coded temporary
  source-package path.  That directory existed at replay time, and its two
  CNFs and two DRUP files were byte-identical to the durable package
  copies.  The script completed successfully, but is not portable as
  written.
\end{itemize}

There is no Lean theorem of \(z(20)=6\) in the bundle.
\artifact{lean_blueprint.md} is a design document that explicitly leaves
the catalogue, residual enumeration, and graph-to-CNF reflection as
formalisation obligations.

\statusbox{The accurate present claim is: \emph{a complete, locally
regenerated candidate computer-assisted proof of \(z(20)=6\), with
certificate-backed SAT refutations and an explicit trusted-computing
boundary}.}

\section{Reproduction record and artifact map}

\subsection{Versions bound to this note}

\begin{itemize}
  \item Local bundle: \artifact{/Users/admin/z20_repro}, Git commit
  \sha{61e8300}.
  \item Public artifact snapshot: \url{https://github.com/ipitchford/z20-cochromatic},
  commit
  \sha{8d9829edebecdfc78384eef5b41617ddd32b11ea}.
  \item Verification date: 26 July 2026, macOS on arm64.
  \item Companion replay receipt:
  \artifact{z20_equals_6_replay_receipt_2026-07-26.md}.
\end{itemize}

\subsection{Load-bearing local files}

\begin{longtable}{@{}>{\raggedright\arraybackslash\footnotesize}p{0.45\textwidth}
>{\raggedright\arraybackslash}p{0.47\textwidth}@{}}
\toprule
Relative path under \texttt{z20\_repro} & Role\\
\midrule
\endhead
\path{nauty/gen44.c} &
Orderly generator for all \((4,4)\)-graphs, using nauty canonical forms.\\
\path{nauty/gen44d.c} &
Dumping variant used to stream every order-\(12\) representative.\\
\path{nauty/cochk.c} &
Exact recursive \(k\)-cocolourability decision procedure.\\
\path{certificate/z20_core_certificate_candidate/generate.py} &
Deterministic fixed-core CNF generator.\\
\path{certificate/z20_core_certificate_candidate/z20_core0.cnf} and
\path{certificate/z20_core_certificate_candidate/z20_core1.cnf} &
The two \(64\)-variable formulas.\\
\path{certificate/z20_core_certificate_candidate/z20_core0.drup} and
\path{certificate/z20_core_certificate_candidate/z20_core1.drup} &
Tree-like RUP/DRUP refutations.\\
\path{verify_certificate.py} &
Separately implemented formula construction, semantic comparison, and RUP
checker; path repair is required as noted above.\\
\path{drat/drat-trim.c} &
Source of the standard DRAT checker and LRAT converter.\\
\path{drat/mine_core0.cnf} and
\path{drat/mine_core1.cnf} &
Independently regenerated, byte-identical CNFs.\\
\path{cake_lpr/cake_lpr_arm8.S} &
Preserved CakeML-compiled checker assembly.\\
\path{cake_lpr/core0.lrat} and
\path{cake_lpr/core1.lrat} &
LRAT proofs accepted by \texttt{cake\_lpr}.\\
\path{nauty/gen44_result.txt},
\path{nauty/z12_result.txt}, and
\path{nauty/lower_bound_result.txt} &
Stored computation transcripts.\\
\bottomrule
\end{longtable}

\subsection{Minimal verdict checks}

The certificate package documents exact build commands.  A reproducer
should, at minimum, check all of the following:

\begin{enumerate}
  \item the package SHA-256 manifest;
  \item exact equality of regenerated and packaged CNFs;
  \item the package RUP checker on both formula/proof pairs;
  \item \texttt{drat-trim} reporting \texttt{s VERIFIED} on both;
  \item \texttt{cake\_lpr} printing the exact verdict
  \texttt{s VERIFIED UNSAT} on both;
  \item a cross-core negative control that is rejected;
  \item the catalogue counts and exact order-\(16\) graph6 strings;
  \item the \(12{,}346\)-representative \(z(8)\) run and the
  \(1{,}449{,}166\)-representative residual \(z(12)\) run.
\end{enumerate}

\section{Provenance and AI disclosure}

The local repository reports that every mathematical step in the result
was produced by AI systems, while the human participant selected the
problem and mediated between systems.  The reported division of labour is:
an unsuccessful Hermes-driven CEGIS search using GPT-5.6 Sol; subsequent
instrumentation and survivor analysis by Claude 4.8; the two-core reduction,
CNF generator, DPLL solver, and DRUP certificates from GPT-5.6 Sol; and
verification by a different AI system, catalogue and residual enumeration,
lower-bound reconstruction, and the earlier write-ups by Claude Opus 5.
This paragraph records repository provenance; it is not an externally
audited authorship determination.

This PDF was prepared with OpenAI Codex from the local proof manuscript,
walkthrough, source code, certificates, logs, and fresh replays listed
above.  Codex revised the exposition, corrected assurance wording, added
the hidden catalogue-pruning lemma, and typeset the document.  The status
label on the title page is intentionally retained.  Any claim that the
problem is settled, or any stronger publication claim, should await
independent human review.

\small\raggedright
\begin{thebibliography}{99}

\bibitem{lesniak-straight}
L. Lesniak-Foster and H. Joseph Straight,
``The cochromatic number of a graph,''
\emph{Ars Combinatoria} \textbf{3} (1977), 39--45.

\bibitem{erdos758}
T. F. Bloom,
``Erd\H{o}s Problem \#758,''
\url{https://www.erdosproblems.com/758},
accessed 26 July 2026.

\bibitem{greenwood-gleason}
R. E. Greenwood and A. M. Gleason,
``Combinatorial relations and chromatic graphs,''
\emph{Canadian Journal of Mathematics} \textbf{7} (1955), 1--7.
\url{https://doi.org/10.4153/CJM-1955-001-4}.

\bibitem{mckay-catalogue}
B. D. McKay,
``Ramsey graphs,'' Combinatorial Data collection,
\url{https://users.cecs.anu.edu.au/~bdm/data/ramsey.html},
accessed 26 July 2026.

\bibitem{mckay-piperno}
B. D. McKay and A. Piperno,
``Practical graph isomorphism, II,''
\emph{Journal of Symbolic Computation} \textbf{60} (2014), 94--112.
\url{https://doi.org/10.1016/j.jsc.2013.09.003}.

\bibitem{drat-trim}
N. Wetzler, M. J. H. Heule, and W. A. Hunt, Jr.,
``DRAT-trim: Efficient checking and trimming using expressive clausal
proofs,'' in \emph{Theory and Applications of Satisfiability Testing --
SAT 2014}, LNCS 8561, 422--429.
\url{https://doi.org/10.1007/978-3-319-09284-3_31}.

\bibitem{cake-lpr}
Y. K. Tan, M. J. H. Heule, and M. O. Myreen,
``Verified propagation redundancy and compositional UNSAT checking in
CakeML,''
\emph{International Journal on Software Tools for Technology Transfer}
\textbf{25} (2023), 167--184.
\url{https://doi.org/10.1007/s10009-022-00690-y}.

\bibitem{z20repo}
\emph{\(z(20)=6\): a candidate computer-assisted proof},
public artifact repository, commit
\sha{8d9829edebecdfc78384eef5b41617ddd32b11ea},
\url{https://github.com/ipitchford/z20-cochromatic},
accessed 26 July 2026.

\end{thebibliography}
\normalsize

\end{document}
